% ===================================================================
%  TABELL Nr. 14 — D'Strooss. --- D'Händler.
%  Traduction 2026 d'apres texto/io, controlee sur texto/fr, texto/de
%  et texto/nl. Consigne : voir texto/lb/10-tabell-01.tex.
%
%  SOIXANTE METIERS EN SEPT PAGES, ET UNE SEULE QUESTION SUFFIT A
%  PREDIRE LE NOM DE CHACUN : le metier avait-il une corporation dans
%  la vieille ville, ou est-il apparu dans la rue commercante du
%  XIXe siecle ?
%
%  Avaient une corporation, et portent un nom germanique :
%  « Bäcker » le boulanger, « Metzler » le boucher, « Schmadd » le
%  forgeron, « Schneider » le tailleur, « Schouster » le cordonnier,
%  « Faassbënner » le tonnelier, « Auermécher » l'horloger,
%  « Buchbënner » le relieur, « Spengler » le ferblantier,
%  « Kaminsfeeër » le ramoneur, « Färber » le teinturier,
%  « Waffenschmadd » l'armurier.
%
%  N'en avaient pas, et portent un nom francais : « Coiffeur »,
%  « Fotograf », « Optiker », « Dentist », « Pâtissier »,
%  « Bijoutier », « Modistin », « Tapisséier »,
%  « Commissionnaire », « Gouvernante », « Affekot »,
%  « Notaire », « Marquise », « Étalage », « Imperial ».
%
%  LE TABLEAU 14 LITUANIEN AVAIT OPPOSE L'ATELIER ET LE DIPLOME ;
%  CELUI-CI OPPOSE LA CORPORATION ET LA VITRINE. La ligne n'est plus
%  celle du savoir mais celle de la DATE : ce qui existait avant que
%  le francais fut la langue des affaires a garde son nom, ce qui est
%  venu apres a pris le sien. C'est le critere de la date, invente
%  pour l'irlandais et verifie sur trente plantes en lituanien,
%  applique ici a une rue entiere.
%
%  ET LE PARTAGE EST DATABLE, ce qui n'etait le cas d'aucun des
%  precedents : on peut ouvrir le registre des corporations de la
%  vieille ville et y lire d'avance quels metiers de ce tableau
%  s'ecrivent en allemand.
% ===================================================================

%%K t14-apar-1 apar
\VUsaut{4.0mm}
\VUtitre{15.0pt}{60}{TABELL Nr. 14}
\VUsaut{3.0mm}

%%K t14-tit-1 sub
\VUpk{11.4pt}{40}{D'Strooss. --- D'Händler.}
\VUsaut{3.0mm}

%%K t14-01-1 p
1. --- Mäi Frënd Jang, dee op dem Land wunnt, ass komm, fir dräi Deeg a menger
Famill ze verbréngen. \VUgras{Bis} elo hat hien nach ni d'Geleeënheet, eng grouss
Stad ze gesinn; dofir verbrénge mir all eis Deeg mat Spadséieren, an ech erkläre him
dat, wat hien net kennt.

%%K t14-02-1 p
2. --- Fir d'éischt si mir d'\VUgras{Observatoire} \textsuperscript{(1)} besiche
gaangen, dat hien ganz interesséiert huet. Wéi mir duerch d'\VUgras{Strooss}
\textsuperscript{(3)} zréckkomm sinn, wou d'\VUgras{Tierkesch Buedanstalt}
\textsuperscript{(2)} ass, hu mir e \VUgras{Begriefnes} \textsuperscript{(4)}
begéint. Virum \VUgras{Leichenzuch} ass de \VUgras{Klerus} gaangen, virum
\VUgras{Leichewon,} an deem de \VUgras{Sarg} louch. Vill Frënn hunn déi
\VUgras{trauereg} Famill begleet.

%%K t14-02-2 p
Nodeems de Cortège laanschtgaangen ass, si mir iwwert d'Strooss gaangen, virun der
grousser \VUgras{Brauerei} \textsuperscript{(5)}, an där vill \VUgras{Konsumenten}
op der Terrass \VUgras{Béier} gedronk hunn, geschützt vun der verglaster
\VUgras{Marquise} \textsuperscript{(6)}.

%%K t14-03-1 p
3. --- De Jang war erstaunt iwwert dat \VUgras{Wopeschëld,} dat tëscht dem Magasin
vum \VUgras{Likörhändler} \textsuperscript{(7)} an deem vum
\VUgras{Museksverkeefer} \textsuperscript{(10)} hänkt. Hien huet mech no senger
Bedeitung gefrot; ech hunn geäntwert, datt et dat englescht \VUgras{Konsulat}
\textsuperscript{(8)} bezeechent, an ech hunn him de \VUgras{Konsul}
\textsuperscript{(9)} gewisen, dee um Balkon ass.

%%K t14-tit-2 sub
\VUpk{10.2pt}{40}{Buttecker kucken.}
\VUsaut{3.0mm}

%%K t14-04-1 p
4. --- Natierlech si mir virun der Vitrine mat de \VUgras{Museksinstrumenter,} den
\VUgras{Harmoniumen,} den \VUgras{Uergelen} an de \VUgras{Pianoen} stoe bliwwen; mir
hunn déi illustréiert Ëmschléi vun de \VUgras{Partituren} a vun de
\VUgras{Museksstécker} fir Piano gekuckt, an hunn d'\VUgras{Lyra} aus
vergoldenem Holz bewonnert, déi als Schëld déngt.

%%K t14-05-1 p
5. --- D'Stëbswolleke vun de \VUgras{Stroossekier} \textsuperscript{(11)} a vum
\VUgras{Kierwon} \textsuperscript{(12)} hunn eis gezwongen, séier virun de schéine
\VUgras{Miwwelen} vum \VUgras{Miwwelhändler} \textsuperscript{(13)}
laanschtzegoen a virun der Vitrine mat \VUgras{Uewen} a \VUgras{Lampen} vum
\VUgras{Spengler} \textsuperscript{(14)}.

%%K t14-06-1 p
6. --- Iwwregens hu mir op där Plaz misse vum Trottoir erof, well en war mat Päck
verstellt, déi ee vun engem \VUgras{Transportwon} \textsuperscript{(16)} geholl huet,
fir se an e \VUgras{Magasin} \textsuperscript{(15)} ze setzen, dee grad gelount ginn
war.

%%K t14-06-2 p
Dat huet eis verhënnert, déi schéi Kutsche vum \VUgras{Kutschebauer}
\textsuperscript{(17)} ze gesinn, mä huet eis net verhënnert, stoen ze bleiwen a de
\VUgras{Packer} \textsuperscript{(19)} ze kucken, deen eng Këscht fir säin Noper
mécht, de \VUgras{Vëlo- an Autoverkeefer} \textsuperscript{(18)}. Duerno, well et
scho spéit war, si mir heemgaangen.

%%K t14-tit-3 sub
\VUpk{10.2pt}{40}{Spadséiergang duerch d'Stad.}
\VUsaut{3.0mm}

%%K t14-07-1 p
7. --- Den Dag drop huet meng Mamm dem Jang proposéiert, sech fotograféieren ze
loossen, duerno huet si mech gebieden, hien bei de \VUgras{Fotograf}
\textsuperscript{(20)} ze begleeden. Nom Mëttegiesse si mir an d'\VUgras{Atelier} vum
Kënschtler gaangen, wou mir zimlech laang hu misse waarden. Fir eis ze beschäftegen,
hu mir d'\VUgras{Portraiten} \textsuperscript{(22)} gekuckt, déi un der Mauer hänken.

%%K t14-07-2 p
Wéi mir un der Rei waren, huet d'Aarbecht net laang gedauert, a soubal mäi Frënd
fäerdeg war, virum \VUgras{Fotoapparat} \textsuperscript{(21)} ze poséieren, si mir
erofgaangen.

%%K t14-08-1 p
8. --- Um éischte Stack hu mir duerch d'Dier vum \VUgras{Coiffeur}
\textsuperscript{(23)}, déi op stoung, säi Gesell gesinn, deen engem Client d'Hoer
geschnidden huet.

%%K t14-08-2 p
Wéi mir op d'Strooss komm sinn, si mir virum \VUgras{Tubaksbuttek}
\textsuperscript{(24)} laanschtgaangen. De Jang war esou ameséiert vun der rouder
\VUgras{Laterne} \textsuperscript{(25)} an der \VUgras{décker Päif,} déi als
\VUgras{Schëld} \textsuperscript{(26)} déngt, datt hien mat zwee alen Hären
zesummegestouss ass, déi geschwat hunn, während si \VUgras{Tubak geschnauft}
\textsuperscript{(27)} hunn.

%%K t14-tit-4 sub
\VUpk{10.2pt}{40}{D'Konditerei.}
\VUsaut{3.0mm}

%%K t14-09-1 p
9. --- D'\VUgras{Bankbilljeeën} a \VUgras{Suestécker} vun alle Länner, déi an der
\VUgras{Vitrine} vum \VUgras{Wiessler} \textsuperscript{(29)} ausgestallt sinn, hunn
eis Opmierksamkeet e Moment ugezunn, mä well et d'Stonn vum Goûter war, si mir bei
de \VUgras{Pâtissier} \textsuperscript{(31)} eragaangen, ouni méi laang stoen ze
bleiwen.

%%K t14-10-1 p
10. --- Wéi mir all d'\VUgras{Séissegkeeten} gesinn hunn, déi de Magasin huet ---
\VUgras{Kichelcher, Hunnegkuch, Biscuiten} an all Zorte \VUgras{Bonbongen} ---, konnte
mir net liicht eppes auswielen. Endlech huet de Jang de \VUgras{Cremschnitt}
gewielt, an ech hunn nëmmen eng kleng \VUgras{Kiischtentaart} geholl, well dat Séisst
\VUgras{mécht mer d'Zänn wéi,} an ech wëll guer net ze dacks bei den
\VUgras{Dentist} \textsuperscript{(30)} goen.

%%K t14-tit-5 sub
\VUpk{10.2pt}{40}{Maschinnschreiwen.}
\VUsaut{3.0mm}

%%K t14-11-1 p
11. --- Fir mengem Frënd d'\VUgras{Schreifmaschinn} ze weisen, vun där hien héieren
hat, déi hien awer net kannt huet, si mir an de \VUgras{Büro} \textsuperscript{(32)}
vu mengem \VUgras{Cousin} \textsuperscript{(34)} gaangen. Mir hunn hien fonnt, wéi hie
seng Bréiwer sengem \VUgras{Sekretär} \textsuperscript{(35)} diktéiert huet, deen e
\VUgras{Stenograf} ass. Kaum war e Bréif fäerdeg, do huet e
\VUgras{Maschinnschreiwer} \textsuperscript{(33)} en op senger Maschinn ofgeschriwwen.
Well mir d'Aarbecht vu mengem Verwandten net wollte stéieren, si mir nëmmen e puer
Ablécker bei him bliwwen.

%%K t14-12-1 p
12. --- Ënner dem Büro vu mengem Cousin wunnt e groussen
\VUgras{Auermécher-Bijoutier} \textsuperscript{(37)}, deen ausserdeem Goldschmadd
ass. Säi herrleche Magasin huet vill \VUgras{Aueren, Penduelen, Poschaueren,}
\VUgras{Kettelcher} a wäertvoll \VUgras{Bijouen,} zum Beispill
\VUgras{Pärelekëtten,} gëllen \VUgras{Armbänder, Ouerréng} a \VUgras{Réng,} déi mat
\VUgras{Edelsteng} an \VUgras{Diamanten} geschmückt sinn. Eng vun de Vitrinne
ass extra fir d'\VUgras{Goldwueren} reservéiert an huet eng wichteg Sammlung vu
sëlwer \VUgras{Bestecker.}

%%K t14-13-1 p
13. --- Wéi ech ëm den Eck vun der Strooss gaange sinn, hunn ech gemierkt, datt ee
d'\VUgras{Plack} \textsuperscript{(36)} gewiesselt huet, déi nieft der
\VUgras{Nummer} vum Haus hänkt. Ech wollt meng Bemierkung grad haart soen, wéi déi
frou Kläng vun \VUgras{der} Militärmusek ze héiere waren. Mir sinn dem Regiment
entgéintgelaf, ouni ze denken, datt \VUgras{et} virun de Magasinen vum
\VUgras{Hutthändler} \textsuperscript{(39)} a vum \VUgras{Handschuenhändler}
\textsuperscript{(40)} laanschtgoe géif, an datt mir grad esou gutt gestanen hätten,
fir säin Defilé ze gesinn, wa mir virun der Vitrine vum \VUgras{Optiker}
\textsuperscript{(38)} bliwwe wieren, déi ganz voll ass mat \VUgras{Pincenez,
Ouerbrëller} a \VUgras{Ferngläser.}

%%K t14-tit-6 sub
\VUpk{10.2pt}{40}{Defilé am Marsch.}
\VUsaut{3.0mm}

%%K t14-14-1 p
14. --- Virum \VUgras{Regiment} \textsuperscript{(41)} ass den
\VUgras{Tamburmajor} \textsuperscript{(42)} marschéiert, gefollegt vun den
\VUgras{Trommler} \textsuperscript{(43)}, den \VUgras{Trompetter}
\textsuperscript{(44)} an der ganzer \VUgras{Musek.} De \VUgras{General}
\textsuperscript{(45)}, deen op der Plaz mat sengem Adjutant war, war nieft dem
\VUgras{Waasserstrahl} \textsuperscript{(49)} vun der \VUgras{Fontän}
\textsuperscript{(48)} stoe bliwwen, fir dem \VUgras{Defilé} bäizewunnen.

%%K t14-14-2 p
D'\VUgras{Stabsoffizéier} \textsuperscript{(46)}, de \VUgras{Colonel} an de
\VUgras{Leitnant-Colonel} hunn hie mam Degen gegréisst. D'\VUgras{Majoren} waren
virun hire \VUgras{Batailloner} \textsuperscript{(47)}, an all \VUgras{Kapitän} huet
seng \VUgras{Kompanie} kommandéiert, während d'\VUgras{Leitnanten,}
\VUgras{Ënnerleitnanten} an \VUgras{Ënneroffizéier} hir gewéinlech Plaz nieft hire
Männer ageholl hunn.

%%K t14-15-1 p
15. --- Vill Passanten hu sech op der \VUgras{Kräizung} \textsuperscript{(50)}
versammelt; d'Händler --- de \VUgras{Prabbelimécher} \textsuperscript{(51)}, de
\VUgras{Färber} \textsuperscript{(53)}, de \VUgras{Waffenschmadd}
\textsuperscript{(54)} --- si erauskomm, fir d'\VUgras{Zaldoten} ze gesinn. All
Gefierer --- \VUgras{Fiakeren} \textsuperscript{(52)}, \VUgras{Härewon}
\textsuperscript{(57)} an och de \VUgras{Spritzewon} \textsuperscript{(56)} --- hu
sech laanscht den Trottoir gestallt.

%%K t14-tit-7 sub
\VUpk{10.2pt}{40}{Zeenen op der Strooss.}
\VUsaut{3.0mm}

%%K t14-15-2 p
En \VUgras{Handelsreesenden} \textsuperscript{(55)}, deen a senger Hand seng
\VUgras{Musterkëschten} gedroen huet, ass séier iwwert d'Strooss gaangen, an
d'Leit, déi um \VUgras{Imperial} \textsuperscript{(59)} vum \VUgras{Omnibus}
\textsuperscript{(58)} souzen, hunn sech ëmgedréint, fir dat Schauspill ze
genéissen. En aarme \VUgras{Stroosse-Sänger} \textsuperscript{(60)} mat senger
\VUgras{Dréiuergel} \textsuperscript{(61)} huet säi \VUgras{Lidd} misse
mattendran ofbriechen, well de Kaméidi vun den Trommelen seng Stëmm iwwerdeckt
huet.

%%K t14-16-1 p
16. --- Wéi d'Regiment virun de \VUgras{Stoffbuttek} \textsuperscript{(64)} komm ass,
hunn d'\VUgras{Näischtesch} \textsuperscript{(63)} an d'\VUgras{Zouschneider}
\textsuperscript{(62)} hir \VUgras{Neimaschinnen} an hir Aarbecht gelooss, fir duerch
d'Fënster ze kucken. De \VUgras{Händler} \textsuperscript{(65)}, deen grad nei
\VUgras{Kostümer} \textsuperscript{(66)} an säin \VUgras{Étalage} gehaangen hat, huet
d'\VUgras{Duch} \textsuperscript{(67)} an d'\VUgras{Stoffer} eranhuele gelooss, déi
um Trottoir nieft dem \VUgras{Kanalisatiounslach} \textsuperscript{(68)}
opgestallt waren, aus Angscht, d'Masse géif se ëmwerfen; souguer en alen
\VUgras{Hauséierer} \textsuperscript{(69)}, mat deem eng Porteusse ëm Strëmp
gemarchandéiert huet, war gezwongen, an e Corridor eranzegoen, fir säi Verkaf fäerdeg
ze maachen.

%%K t14-16-2 p
Mir hunn d'Regiment bis bei d'Kasär begleet, \VUgras{an,} ganz zefridde mat eisem
Dag, si mir zur Stonn vum Owesiesse zréckkomm.

%%K t14-tit-8 sub
\VUpk{10.2pt}{40}{Verschidde Beruffer a Métieren.}
\VUsaut{3.0mm}

%%K t14-17-1 p
17. --- Den Dag drop huet mäi Papp eis proposéiert, d'Dréckerei vun der lokaler
Zeitung ze besichen. Mir hunn et mat Begeeschterung ugeholl.

%%K t14-17-2 p
Den \VUgras{Drécker} ass och \VUgras{Buchhändler} \textsuperscript{(70)}
\VUgras{(= Bicherverkeefer),} \VUgras{Verleeër,} \VUgras{Pabeierhändler} a
\VUgras{Buchbënner}. Um éischte Stack huet hien d'\VUgras{Redaktioun}
\textsuperscript{(71)} vun der Zeitung ageriicht, an och d'\VUgras{Gravéieratelier,}
an deem d'\VUgras{Lithografen} \textsuperscript{(72)} an d'\VUgras{Kopferstiecher}
\textsuperscript{(73)} schaffen, a schlussendlech d'\VUgras{Setzerei,} an där
d'\VUgras{Dréckereiaarbechter} \textsuperscript{(74)} sinn. Déi eigentlech
\VUgras{Dréckerei} \textsuperscript{(75)}, d'\VUgras{Dréckerpressen} an déi
verschidde Maschinne sinn am Parterre.

%%K t14-tit-9 sub
\VUpk{10.2pt}{40}{De Jang stellt ganz vill Froen.}
\VUsaut{3.0mm}

%%K t14-18-1 p
18. --- Wéi mir aus der Dréckerei erausgaange sinn, ass de Jang ganz erstaunt virun
enger klenger verglaster Këscht stoe bliwwen, déi op enger Kolonn nieft dem
\VUgras{Kuerzwuerebuttek} \textsuperscript{(77)} ass, an huet gefrot, wat se soll.
Mäi Papp huet him erkläert, datt dat e \VUgras{Feiermelder} \textsuperscript{(76)}
ass. Iwwregens war alles an der Stad fir mäi Frënd erstaunlech, vun der einfacher
\VUgras{Steefontän} \textsuperscript{(78)} an dem liichte \VUgras{Lastdräirad}
\textsuperscript{(80)}, dat e \VUgras{Groom} \textsuperscript{(79)} a Livrée féiert,
bis bei de \VUgras{Postwon} \textsuperscript{(81)}.

%%K t14-19-1 p
19. --- Während mäi Papp eis e puer Ablécker verlooss huet, fir mam
\VUgras{Steiereinnehmer} \textsuperscript{(83)} ze schwätzen, deen ugehalen hat, fir
sech mat engem \VUgras{Affekot} \textsuperscript{(82)} an engem \VUgras{Notaire}
\textsuperscript{(84)} z'ënnerhalen, hu mir eis ameséiert, andeem mir mat den Ae
de Verkéier op der Strooss verfollegt hunn. Déi verschiddenste Leit sinn viru eis
laanschtgaangen: e \VUgras{Pâtissiersgesell} \textsuperscript{(85)}, deen an engem
Kuerf eng herrlech \VUgras{Pastéit} gedroen huet, ass hannert engem
\VUgras{Kleederverkeefer} \textsuperscript{(86)} komm; e \VUgras{Laterneunzënner}
\textsuperscript{(87)} huet mat engem \VUgras{Gasmann} \textsuperscript{(88)}
geschwat; eng \VUgras{Klouschterfra} \textsuperscript{(89)}, déi e Weesemeedche bei
der Hand gehalen huet, ass an hiert Klouschter zréckgaangen.

%%K t14-tit-10 sub
\VUpk{10.2pt}{40}{Op eisem Heemwee.}
\VUsaut{3.0mm}

%%K t14-20-1 p
20. --- Am Ablack, wou mäi Papp eis nokomm ass, huet de Jang gelaacht, wéi hien
dat dreckegt Gesiicht an déi komesch Kleedung vun zwee ganz roussege
\VUgras{Kaminsfeeër} \textsuperscript{(91)} gesinn huet.

%%K t14-21-1 p
21. --- Ech wollt vun der \VUgras{Blummefra} \textsuperscript{(92)} eng Planz fir
meng Mamm kafen, mä mäi Papp huet mer gewisen, datt se ganz schwéier ze droe wier an
datt et besser wier, \VUgras{Orangen} ze kafen. Mir sinn no bei d'\VUgras{Verkeeferin}
\textsuperscript{(94)} gaangen, a wéi d'\VUgras{Gouvernante} \textsuperscript{(93)},
déi fir hiren Zögling eppes eraus gesicht huet, hire Kaf fäerdeg hat, hu mir eis
bedéngen gelooss.

%%K t14-22-1 p
22. --- Op dem Heemwee hu mir e puer Bekannter begéint: eng al Frëndin vu menger
Famill, déi sech um Aarm vun hirer \VUgras{Begleeterin} \textsuperscript{(95)}
gestäipt huet, den \VUgras{Ingenieur} \textsuperscript{(96)} vun de Brécken a
Weeër, an eng vu menge Cousinne mat hirer \VUgras{Kannerfra}
\textsuperscript{(98)}.

%%K t14-22-2 p
Duerno hu mir d'\VUgras{Modistin} \textsuperscript{(97)} vu menger Mamm begéint an
zwee \VUgras{Mönche} \textsuperscript{(99)}, déi der Stad friem waren an déi bei eis
komm sinn, fir mäi Papp no dem Wee bei d'Gare ze froen.

\VUsaut{6.0mm}
\VUfilet{22mm}
